\section{Conclusion}
\label{sec:conclusion}

We showed that compound sparsity severity is not phase-invariant but exhibits an emergent phase-strategy interaction. Under importance-based eviction, the paired median difference $\Delta$ increases from $0.39$ to $8.48$ (DeepSeek, $n{=}496$; Spearman $\rho{=}1.0$) and $0.07$ to $3.16$ (Qwen, $n{=}475$), all $p < 10^{-13}$. Random eviction is $2.7$--$17{\times}$ weaker where both strategies are significant ($s \geq 0.25$; DeepSeek $15$--$17{\times}$, Qwen $2.7$--$3.5{\times}$). The importance phase gradient replicates on LlamaForCausalLM ($d_z{=}0.81$--$1.15$, $n{=}250$, all $p < 10^{-29}$), supporting architecture-generality of phase modulation, while strategy specificity is architecture-dependent ($2.0$--$2.6{\times}$ on LLaMA vs.\ $15$--$17{\times}$ on DeepSeek). A model-dependent crossover on DeepSeek demonstrates that no single strategy dominates uniformly. Per-phase Gini analysis reveals importance discrimination collapse as a correlational diagnostic: despite near-identical aggregate Gini drops, severity diverges by $2.7$--$17{\times}$, implicating the structural specificity of importance-based pruning.
Beyond the empirical findings, the phase-gated analysis framework itself---decomposing compound severity by reasoning phase and eviction strategy---transforms a scalar interaction metric into a structured diagnostic that reveals where compound sparsity fails and why. This framework provides practitioners with a principled basis for phase-aware sparsity allocation: applying aggressive compression during high-sparsity (structured computation) phases while preserving budget during low-sparsity (exploratory reasoning) phases. Future work could explore generalization to other reasoning domains, causal validation of importance discrimination collapse, and systematic characterization of asymmetric compound interactions.
